\begin{exblock}{复习题 10}

1. 设 $f$ 为单变量函数, 连续可导. 令
\[
F(t) = \iint\limits_{[0,t]^{2}} f(xy) \mathrm{d}x\mathrm{d}y.
\]
证明:
\[
F'(t) = \frac{2}{t} \left( F(t) + \iint\limits_{[0,t]^{2}} xy f'(xy) \mathrm{d}x\mathrm{d}y \right).
\]

2. 设 $f$ 为单变量的连续函数. 令
\[
F(t) = \iint\limits_{[0,t]^{2}} f(xy) \mathrm{d}x\mathrm{d}y.
\]
证明:
\[
F'(t) = \frac{2}{t} \int_{0}^{t^{2}} f(u) \mathrm{d}u.
\]

3. 设 $f$ 为单变量函数, 连续可导. 令
\[
F(t) = \iiint\limits_{[0,t]^{3}} f(xyz) \mathrm{d}x\mathrm{d}y\mathrm{d}z.
\]
证明:
\[
F'(t) = \frac{3}{t} \left( F(t) + \iiint\limits_{[0,t]^{3}} xyz f'(xyz) \mathrm{d}x\mathrm{d}y\mathrm{d}z \right).
\]
能否将题 1 与题 3 的结果推广到一般情形.

4. 设 $f$ 为单变量的连续函数. 令
\[
F(t) = \iiint\limits_{[0,t]^{3}} f(xyz) \mathrm{d}x\mathrm{d}y\mathrm{d}z.
\]
证明:
\[
F'(t) = \frac{3}{t} \int_{0}^{t^{3}} \frac{\displaystyle\int_{0}^{u} f(s) \mathrm{d}s}{u} \mathrm{d}u.
\]

5. 设 $\Omega$ 是由在第一象限里的圆周 $x^{2}+y^{2}=1$ 与直线 $x+y=1$ 所围成的图形. 证明:
\[
\iint\limits_{\Omega} \frac{\mathrm{d}x\mathrm{d}y}{xy\left( \ln^{2}x + \ln^{2}y \right)}
= \frac{\pi}{2} \ln 2.
\]

6. 设平面区域 $\Omega$ 在 $x$ 轴与 $y$ 轴的投影长度分别为 $l_{x}$ 与 $l_{y}$,
$\Omega$ 的面积为 $v(\Omega)$, $(\alpha, \beta)$ 为 $\Omega$ 内任一点. 证明:
\begin{enumerate}
\item[(1)] $\displaystyle \left| \iint\limits_{D} (x-\alpha)(y-\beta) \mathrm{d}x\mathrm{d}y \right|
\leqslant l_{x} l_{y} v(\Omega)$;
\item[(2)] $\displaystyle \left| \iint\limits_{D} (x-a)(y-\beta) \mathrm{d}x\mathrm{d}y \right|
\leqslant \frac{1}{4} l_{x}^{2} l_{y}^{2}$
\end{enumerate}

7. 设
\[
f(x,y) = \begin{cases}
\dfrac{1}{q_{x}} + \dfrac{1}{q_{y}}, & \text{当 } (x,y) \text{ 为 } [0,1]^{2} \text{ 中有理点},\\
0, & \text{当 } (x,y) \text{ 为 } [0,1]^{2} \text{ 中非有理点},
\end{cases}
\]
其中 $q_{x}$ 表示有理数 $x$ 化成既约分数后的分母.
证明: $f(x,y)$ 在 $[0,1]^{2}$ 上的二重积分存在而两个累次积分不存在.

8.（Pringsheim 定理）设
\[
f(x,y) = \begin{cases}
1, & \text{当 } (x,y) \text{ 为 } [0,1]^{2} \text{ 中有理点, 且 } q_{x} \neq q_{y} \text{ 时},\\
0, & \text{当 } (x,y) \text{ 为 } [0,1]^{2} \text{ 中其他点时},
\end{cases}
\]
其中 $q_{x}$ 表示有理数 $x$ 化成既约分数后的分母.
证明: $f(x,y)$ 在 $[0,1]^{2}$ 上的二重积分不存在而两个累次积分都存在且为 0.

9. 证明: 由曲面
$\dfrac{x^{n}}{a^{n}} + \dfrac{y^{n}}{b^{n}} + \dfrac{z^{n}}{c^{n}} = 1$,
$x=0, y=0, z=0\ (n > 0, a > 0, b > 0, c > 0)$
所界立体 $\Omega$ 的体积
\[
v(\Omega) = \frac{abc}{3n^{2}} \frac{\Gamma^{3}\left( \frac{1}{n} \right)}{\Gamma\left( \frac{3}{n} \right)}
\]

10. 证明: 由曲面
$\left( \dfrac{x}{a} + \dfrac{y}{b} \right)^{n} + \left( \dfrac{z}{c} \right)^{m} = 1$,
$x=0, y=0, z=0\ (n > 0, m > 0, a > 0, b > 0, c > 0)$
所界立体 $\Omega$ 的体积
\[
v(\Omega) = \frac{abc}{2m+n}
\frac{\Gamma\left( \frac{1}{m} \right)\Gamma\left( \frac{2}{n} \right)}
{\Gamma\left( \frac{1}{m}+\frac{2}{n} \right)}
\]

11. 设 $\Omega$ 为第一卦限内且由曲面
$z = \dfrac{x^{2}+y^{2}}{m}$,\quad $z = \dfrac{x^{2}+y^{2}}{n}$,\quad
$xy = a^{2}$, $xy = b^{2}$, $y = \alpha x$, $y = \beta x$
\ $(0 < a < b,\ 0 < \alpha < \beta,\ 0 < m < n)$
所围成的立体. 证明:
\[
\iiint\limits_{\Omega} xyz \mathrm{d}x\mathrm{d}y\mathrm{d}z
= \frac{1}{32}\left( b^{8} - a^{8} \right)
\left( \frac{1}{m^{2}} - \frac{1}{n^{2}} \right)
\left[ \left( \beta^{2}-a^{2} \right)
\left( 1 + \frac{1}{\alpha^{2}\beta^{2}} \right) + 4 \ln \frac{\beta}{\alpha} \right]
\]

12. 应用数学归纳法证明: Dirichlet 公式
\begin{align*}
&\idotsint _{x_{1},x_{2},\cdots,x_{n} \geqslant 0}
x_{1}^{p_{1}-1} x_{2}^{p_{2}-1} \cdots x_{n}^{p_{n}-1}
\mathrm{d}x_{1} \mathrm{d}x_{2} \cdots \mathrm{d}x_{n} \notag\\
= \frac{\Gamma(p_{1})\Gamma(p_{2}) \cdots \Gamma(p_{n})}
{\Gamma(p_{1}+p_{2}+\cdots+p_{n+1})}
\notag\\
&\qquad (p_{1}, p_{2}, \cdots, p_{n} > 0)
\end{align*}

13. 设 $f(x)$ 为连续函数, 应用数学归纳法证明: Liouville 公式
\[
\begin{aligned}
&\idotsint _{x_{1},x_{2},\cdots,x_{n} \geqslant 0,\, x_{1}+x_{2}+\cdots+x_{n} \leqslant 1}
f(x_{1}+x_{2}+\cdots+x_{n}) x_{1}^{p_{1}-1} x_{2}^{p_{2}-1} \cdots x_{n}^{p_{n}-1}
\mathrm{d}x_{1} \mathrm{d}x_{2} \cdots \mathrm{d}x_{n} \notag\\
&\quad= \frac{\Gamma(p_{1})\Gamma(p_{2}) \cdots \Gamma(p_{n})}{\Gamma(p_{1}+p_{2}+\cdots+p_{n})}
\int_{0}^{1} f(u) u^{p_{1}+p_{2}+\cdots+p_{n-1}} \mathrm{d}u
\qquad (p_{1}, p_{2}, \cdots, p_{n} > 0) \notag
\end{aligned}
\]

14. 设 $\sum\limits_{i,j=1}^{n} a_{ij} x_{i} x_{j}\ (a_{ij} = a_{ji})$ 为正定二次型. 证明:
\[
\int_{\mathbb{R}^{n}} \cdots \int_{\mathbb{R}^{n}}
\mathrm{e}^{-\left( \sum\limits_{i,j=1}^{n} a_{ij} x_{i} x_{j}
+ 2 \sum\limits_{i=1}^{n} b_{i} x_{i} + c \right)}
\mathrm{d}x_{1} \mathrm{d}x_{2} \cdots \mathrm{d}x_{n}
= \frac{\pi^{\frac{n}{2}}}{\sqrt{\delta_{n}}}
\mathrm{e}^{-\frac{\Delta_{n}}{\delta_{n}}},
\]
其中
\[
\delta_{n} = \left| \begin{matrix}
a_{11} & \cdots & a_{1n}\\ \vdots & & \vdots\\ a_{n1} & \cdots & a_{nn}
\end{matrix} \right|,
\qquad
\Delta_{n} = \left| \begin{matrix}
a_{11} & \cdots & a_{1n} & b_{1}\\ \vdots & & \vdots & \vdots\\
a_{n1} & \cdots & a_{nn} & b_{n}\\ b_{1} & \cdots & b_{n} & c
\end{matrix} \right|
\]

15. 设 $f(x, y)$ 是在 $I = [0, 1]$ 上有定义的实值函数.
对任意的 $(a, b) \in I$, 设 $I(a, b)$ 是以 $(a, b)$ 为中心并且完全含在 $I$ 内的
且各边与 $I$ 的边平行的正方形中的最大的一个.
如果对 $\forall (a, b) \in I$, 有
\[
\iint\limits_{I(x,b)} f(x,y) \mathrm{d}x\mathrm{d}y = 0.
\]
问: $f(x,y)$ 在 $I$ 上恒等于零吗？（1977, 第 38 届美国 Patnam 数学竞赛题）.

16. 设四次连续可导函数 $f(x, y)$ 在 $[0, 1]^{2}$ 的边界上为零, 并且
\[
\left| \frac{\partial^{4} f}{\partial x^{2} \partial y^{2}} \right| \leqslant M.
\]
证明:
\[
\left| \iint\limits_{[0,1]^{2}} f(x,y) \mathrm{d}x\mathrm{d}y \right| \leqslant \frac{1}{144} M.
\]

17. 设二元函数 $f(x, y)$ 在 $[a, b] \times [c, d]$ 上有定义,
并且 $f(x, y)$ 对于确定的 $x \in [a, b]$ 是 $y$ 在 $[c, d]$ 上的增函数,
对于确定的 $y \in [c, d]$ 是 $x$ 在 $[a, b]$ 上的增函数.
证明: $f(x, y)$ 在 $[a, b] \times [c, d]$ 上可积.

18. 设 $f(x,y) \geqslant 0$ 且在 $\left\{ (x,y) \mid x^{2}+y^{2} \leqslant a^{2} \right\}$
上有连续的一阶偏导数, 边界上取值为零. 证明:
\[
\left| \iint\limits_{x^{2}+y^{2} \leqslant a^{2}} f(x,y) \mathrm{d}x\mathrm{d}y \right|
\leqslant \frac{1}{3} \pi a^{3}
\max_{x^{2}+y^{2} \leqslant a^{2}}
\sqrt{ \left( \frac{\partial f}{\partial x} \right)^{2}
+ \left( \frac{\partial f}{\partial y} \right)^{2} }
\]

19. 设 $\Omega = \{ x,y,z,u \mid x,y,z,u \geqslant 0,\ x^{2}+y^{2}+z^{2}+u^{2} \leqslant 1 \}$.
应用

\begin{tabular}{@{}p{0.46\linewidth}p{0.46\linewidth}@{}}
（1）球坐标变换 & （2）双极坐标变换 \\[2pt]
$\displaystyle \begin{cases}
x = r\cos\theta,\\
y = r\sin\theta\cos\theta_{2},\\
z = r\sin\theta\sin\theta_{2}\cos\theta_{3},\\
u = r\sin\theta\sin\theta_{2}\sin\theta_{3};
\end{cases}$
&
$\displaystyle \begin{cases}
x = r\cos\theta,\\
y = r\sin\theta,\\
z = \zeta\cos\varphi,\\
u = \zeta\sin\varphi.
\end{cases}$
\end{tabular}

证明:
\[
\iiint\limits_{\Omega}
\sqrt{\frac{1-x^{2}-y^{2}-z^{2}-u^{2}}{1+x^{2}+y^{2}+z^{2}+u^{2}}}
\mathrm{d}x\mathrm{d}y\mathrm{d}z\mathrm{d}u
= \frac{\pi^{2}}{16} \left( 1 - \frac{\pi}{4} \right)
\]

20. 设 $f_{1}(x), f_{2}(x), \cdots, f_{m}(x)$;
$g_{1}(x), g_{2}(x), \cdots, g_{m}(x)$ 在 $[a,b]$ 上 Riemann 可积. 证明:
\[
\begin{aligned}
&\frac{1}{m!} \int_{a}^{b} \cdots \int_{a}^{b}
\left| \begin{matrix}
f_{1}(x_{1}) & \cdots & f_{1}(x_{m})\\ \vdots & & \vdots\\ f_{m}(x_{1}) & \cdots & f_{m}(x_{m})
\end{matrix} \right|
\left| \begin{matrix}
g_{1}(x_{1}) & \cdots & g_{1}(x_{m})\\ \vdots & & \vdots\\ g_{m}(x_{1}) & \cdots & g_{m}(x_{m})
\end{matrix} \right|
\mathrm{d}x_{1} \mathrm{d}x_{2} \cdots \mathrm{d}x_{m} \notag\\
&\quad= \left| \begin{matrix}
\displaystyle\int_{a}^{b} f_{1}(x) g_{1}(x) \mathrm{d}x & \cdots &
\displaystyle\int_{a}^{b} f_{1}(x) g_{m}(x) \mathrm{d}x\\ \vdots & & \vdots\\
\displaystyle\int_{a}^{b} f_{m}(x) g_{1}(x) \mathrm{d}x & \cdots &
\displaystyle\int_{a}^{b} f_{m}(x) g_{m}(x) \mathrm{d}x
\end{matrix} \right| \notag
\end{aligned}
\]

\end{exblock}
